\documentclass[11pt,a4paper]{article}

\usepackage[utf8]{inputenc}

\usepackage[T1]{fontenc}

\usepackage[polish]{babel}

\usepackage{lmodern}

\usepackage{geometry}

\geometry{margin=27mm}

\usepackage[hidelinks]{hyperref}

\title{Komputacyjny benchmark jakości argumentów i retoryki LLM}

\author{Mateusz Pomianek}

\date{Szkic naukowy, wersja 1.0.0 | 28 listopada 2026}

\begin{document}

\maketitle

\begin{abstract}

Proponowana miara główna to różnica korelacji rang z ocenami ekspertów między dwoma predyktorami.

\end{abstract}

\section{Wprowadzenie}

Ocena argumentacji obejmuje wymiary z teorii i kryteria używane w praktyce (Wachsmuth et al., 2017). Zadania predykcji przekonującego stylu dostarczają odrębnego punktu odniesienia (Habernal \& Gurevych, 2016).

\section{Cel i pytania badawcze}

Hipotezy robocze wymagają zatwierdzenia przed rozpoczęciem właściwej weryfikacji.

\paragraph{H1.} Predyktor wykorzystujący weryfikowalność przesłanek i adekwatność dowodu osiągnie wyższą korelację rang z zaślepioną oceną ekspertów niż predyktor oparty na cechach stylu, w tematach wyłączonych z dopasowania.

\section{Materiał i metody}

Pierwszorzędowy wskaźnik: różnica korelacji rang z ocenami ekspertów między dwoma predyktorami.

Eksperci osobno zakodują jakość dowodu i perswazyjność. Ocena obejmie tę samą pulę odpowiedzi dla obu predyktorów i bootstrap po tematach; rodziny parafraz nie mogą przekraczać granicy rozdziału zbiorów. Wrażliwość zbada podział według modeli generujących.

Przed badaniem z udziałem ludzi należy zatwierdzić wymagania etyczne, zasady zgody i ochronę danych.

\section{Plan analizy}

Zamrożenie estymandu, reguł oceny jakości i obsługi braków nastąpi przed dostępem do danych wynikowych. Raport obejmie oszacowanie kontrastu głównego z przedziałem niepewności i pełne mianowniki. Jeśli przewidziano hipotezy drugorzędowe, przed analizą należy ustalić sposób kontroli wielokrotności.

\section{Wyniki}

W niniejszym pliku nie umieszczono wyników empirycznych, statystyk ani nieprzeprowadzonych analiz. Sekcja wymaga uzupełnienia po uzyskaniu i weryfikacji materiału dowodowego.

\section{Dyskusja i ograniczenia}

Ograniczenia obejmują możliwe błędy pomiaru, dobór próby, alternatywne mechanizmy i ograniczoną możliwość uogólniania. Korelacja nie dowodzi przyczynowości; poprawność techniczna nie jest dowodem znaczenia klinicznego.

\section{Wnioski}

Hipotezy i sposób pomiaru stanowią podstawę oceny wykonalności oraz przygotowania właściwego protokołu.

\section*{Bibliografia}

{\setlength{\parindent}{0pt}
\setlength{\parskip}{5pt}
\hangindent=1.27cm
Habernal, I., \& Gurevych, I. (2016). Which argument is more convincing? Analyzing and predicting convincingness of Web arguments using bidirectional LSTM. \textit{Proceedings of the 54th Annual Meeting of the Association for Computational Linguistics (Volume 1: Long Papers)}, 1589--1599. \url{https://doi.org/10.18653/v1/P16-1150}\par
\hangindent=1.27cm
Wachsmuth, H., Naderi, N., Habernal, I., Hou, Y., Hirst, G., Gurevych, I., \& Stein, B. (2017). Argumentation quality assessment: Theory vs. practice. \textit{Proceedings of the 55th Annual Meeting of the Association for Computational Linguistics (Volume 2: Short Papers)}, 250--255. \url{https://doi.org/10.18653/v1/P17-2039}\par
}

\end{document}
